\documentclass[10pt,a4paper]{amsart}
\usepackage[margin=19mm]{geometry}
\usepackage{amsmath,amssymb,mathtools}
\usepackage[scr=rsfs,cal=euler]{mathalfa}
\usepackage{xfrac}
\usepackage[unicode,hidelinks,bookmarks=false]{hyperref}
\newcommand{\ie}{i.e.\ }
\newcommand{\cf}{cf.\ }
\newcommand{\resp}{respectively\ }
\newcommand{\Subsection}{\subsection}
\providecommand{\nobreakdash}{\nobreak\hbox{-}}
\setlength{\emergencystretch}{2em}
\multlinegap0pt
\usepackage{amssymb}
\usepackage{xcolor}
\newtheorem{lem}{Lemma}
\newcommand{\card}{\operatorname{card}}
\title{The Square Root Problem and Subnormal Aluthge Transforms of Recursively Generated Weighted Shifts}
\author{Raul E. Curto, Hamza El Azhar, Youssef Omari, El Hassan Zerouali}
\date{Source: arXiv:2310.13887v2 (CC BY 4.0)}
\begin{document}
\maketitle

\noindent\emph{Source and licence.} The Square Root Problem and Subnormal Aluthge Transforms of Recursively Generated Weighted Shifts. Version arXiv:2310.13887v2 (CC BY 4.0), distributed by arXiv under the Creative Commons Attribution 4.0 International licence, http://creativecommons.org/licenses/by/4.0/, as recorded in the arXiv licence metadata for this exact version and shown on the abstract page https://arxiv.org/abs/2310.13887v2. Full licence notice: \texttt{licenses/arxiv-2310.13887-CC-BY-4.0.txt}.\par\smallskip

\noindent\textit{Editorial scope.} Counted unit: the article's lem p4-5 (source0.tex lines 422--426). The article's complete original proof of it is delivered with it (source0.tex lines 428--453, one \texttt{proof} environment of 26 lines). No mathematics is written, repaired or re-derived: the statement and the proof are the article's own lines, in the article's own notation, and no retained line is substituted or re-flowed.  Retained uncounted as the source-local context the proof needs: lines 413--415.  Nothing was omitted from any retained span, so the package carries no omission record.  No retained line contains a citation, so the wrapper carries no bibliography and no bibliography entry.  No retained line contains a figure environment, an image-inclusion command or drawing code, so no image is delivered and the package declares no asset dependency.  Editorial formatting only, in addition: an AMS \texttt{amsart} preamble that restates the article's own declarations and macros used by the retained lines, namely the environments \texttt{lem} on separate counters, together with the standard packages those lines need.  The retained environments print in this document as Lemma 1 in order of appearance, whereas the article prints the same items with the numbering of its own full document.  The complete native source of the article as distributed in the arXiv e-print for this exact version is bundled unchanged as \texttt{sources/148830/source0.tex}.
\begin{lem}\label{card3}
    Let $\nu=\sum_{k=0}^q a_k\delta_{\lambda_k}$ be a charge such that $\nu*\nu=\sum_{k=0}^p b_k \delta_{\gamma_k}$ and $\nu*t\nu$ are positive measures. \ If  $\card\Gamma(\lambda^2_k)\leq 3$ for some  $k$, then  $b_k\ne0$.
\end{lem}

\begin{lem}\label{p4-5}
    Let $\nu=\sum_{k=0}^q a_k\delta_{\lambda_k}$ be a charge such that $\nu*\nu=\sum_{k=0}^p b_k \delta_{\gamma_k}$ and $\nu*t\nu$ are positive measures. \ Then \newline
    (i) \ If $q \ge 4$, then  $p\ge 6$. \newline
    (ii) \ If $q \ge 5$, then  $p\ge 7$.
\end{lem}

\begin{proof}
(i) \ Suppose  $q\ge 4$. \ Since  $\card\Gamma(\lambda^2_2)\le 3$ and $\card\Gamma(\lambda_{q-1}^2)\le 3$, we obtain 
$\{\lambda_1^2,\lambda_1\lambda_2,\lambda_2^2, \lambda_{q-1}^2,\lambda_{q-1}\lambda_q,\lambda_q^2\} \subset supp(\mu) $ and hence $p\ge 6$.

(ii) \ From the previous item, $p\ge 6$. \ To show that  $p\ge 7$, it suffices to exhibit a new atom.\\
First, if  $\lambda_2^2$ is $\mathcal{UR}$ or if $\lambda_2^2=\lambda_1\lambda_k$ with some $k>4$, $\lambda_1\lambda_3$  becomes an  $\mathcal{UR}$ and hence provides an additional atom in $\nu*\nu$. \ We write then
  $\lambda_2^2=\lambda_1\lambda_3$ and we show that either $\lambda_3^2$ or $\lambda_2\lambda_3$ is the additional atom or produce a new one. \ To this goal, we suppose that neither $\lambda_3^2$ nor  $\lambda_2\lambda_3$ is $\mathcal{UR}$. \ In this case, necessarily $\lambda_2\lambda_3=\lambda_1\lambda_4$ (otherwise   $\lambda_1\lambda_4$ will be the new atom as an $\mathcal{UR}$.)
We write $\lambda_3^2=\lambda_1\lambda_k$, $\lambda_3^2=\lambda_2\lambda_l$ or  $\lambda_3^2=\lambda_1\lambda_k=\lambda_2\lambda_l$ with zero as  corresponding coefficient. \ Since the two first situations will provide a new atom because of Lemma \ref{card3}, we can assume  that $\lambda_3^2=\lambda_1\lambda_k=\lambda_2\lambda_l$. \ Now, multiplying $\lambda_2\lambda_3=\lambda_1\lambda_4$   with $\lambda_3$ gives $l=4$. 

Now, from the identity $\lambda_1\lambda_k=\lambda_2\lambda_l$, we derive that $k\ge 5$.\\
$1)$ In the case where $k>5$ and  $\lambda_1\lambda_5=\lambda_3\lambda_4$, then by multiplying with $\lambda_3$, we get  
$ \lambda_1\lambda_3\lambda_5=\lambda_3^2\lambda_4=\lambda_1\lambda_k\lambda_4$. \ It follows that  $ \lambda_3\lambda_5=\lambda_k\lambda_4$ for some  $k> 5$, which is impossible. \ Then if $k\ne 5$ $\lambda_1\lambda_5$ will give additional atom as an $\mathcal{UR}$ element. \\
$2)$ $k=5$. \ That is $\lambda_1\lambda_4=\lambda_2\lambda_3$, and
$\lambda_3^2=\lambda_1\lambda_5=\lambda_2\lambda_4$. \ For $r=\frac{\lambda_2}{\lambda_1}$, we get $\lambda_k=\lambda_1r^{k-1}$ for every $k\le 5$. \ Now, to provide the $7^{th}$ atom, it suffices to show that either    $ a_2a_3+a_1a_4\ne 0$ or $a_3^2+2a_1a_5+2a_2a_4\ne 0$. \ Seeking a contradiction, suppose that  $ a_2a_3+a_1a_4=a_3^2+2a_1a_5+2a_2a_4= 0$. \ From the inequality $ (r+r^2)a_2a_3+(1+r^3) a_1a_4\ge 0$, we derive that $a_2a_3<0$. \ Otherwise,
$$ 0\le (r+r^2)a_2a_3+(1+r^3) a_1a_4< (1+r^3)(a_2a_3+ a_1a_4 )= 0$$
It follows also that $a_1a_4>0 $ and then   $a_2a_4>0.$
Now, from 
$a_3^2+2a_1a_5+2a_2a_4=0$, we derive that $a_1a_5<a_1a_5+a_2a_4=-a_3^2<0$ and since 
$(1-r^2)^2>r(1-r)^2$ we obtain the next contradiction
$$\begin{array}{lll} 0 & \le & r^2a_3^2+(1+r^4)a_1a_5+(r+r^3)a_2a_4 \\
&=&-r^2(2a_1a_5+2a_2a_4) +(1+r^4)a_1a_5+(r+r^3)a_2a_4 \\
& = &    (1-r^2)^2a_1a_5+r(1-r)^2a_2a_4 \\
&\le & r(1-r)^2(a_1a_5+a_2a_4) < 0.
\end{array}$$
The proof is complete.
\end{proof}

\end{document}
